\chapter{نتیجه‌گیری}
\label{chap:conclusion}

در این پروژه، مسئله بهینه‌سازی حافظه نهان KV در مدل‌های زبانی بزرگ از دیدگاه ریاضی و الگوریتمی بررسی شد. مهم‌ترین دستاوردها عبارتند از:

\begin{itemize}
\item \textbf{تحلیل ریاضی مکانیزم توجه:} نشان داده شد که پیچیدگی حافظه‌ای توجه استاندارد $O(n^2)$ است و حافظه نهان KV با طول زمینه به‌صورت خطی $O(n)$ رشد می‌کند. این رشد برای مدل‌های بزرگ با طول زمینه بالا به یک تنگنای جدی تبدیل می‌شود.
\item \textbf{مقایسه الگوریتم‌های بهینه‌سازی:} روش‌های GQA، MQA، PagedAttention و FlashAttention از نظر پیچیدگی حافظه، سرعت و دقت مقایسه شدند. نتایج نشان داد که:
  \begin{itemize}
  \item روش GQA با $g=8$ حدود $87.5\%$ در حافظه صرفه‌جویی می‌کند با کاهش دقت کمتر از $0.5\%$.
  \item روش PagedAttention با حذف تکه‌تکه شدن حافظه\footnote{\lr{Fragmentation}}، بهره‌وری حافظه را به $\eta \to 1$ نزدیک می‌کند.
  \item روش FlashAttention با استفاده از کاشی‌کاری\footnote{\lr{Tiling}}، حافظه مصرفی را از $O(n^2)$ به $O(n)$ کاهش می‌دهد.
  \end{itemize}
\item \textbf{شبه‌کد و فلوچارت:} الگوریتم مدیریت حافظه نهان KV با تخصیص صفحه‌بندی شده به‌صورت شبه‌کد ارائه و فلوچارت آن رسم شد.
\end{itemize}

\textbf{پیشنهادات:}
برای ادامه کار، می‌توان روش‌های ترکیبی مانند استفاده هم‌زمان از GQA و FlashAttention را بررسی کرد. همچنین روش‌های جدید مانند سیاست‌های تخلیه حافظه نهان\footnote{\lr{Cache Eviction Policies}} و \lr{StreamingLLM} \cite{xiao2023streamingllm} که امکان پردازش دنباله‌های بی‌نهایت را فراهم می‌کنند، می‌توانند موضوع تحقیقات آینده باشند. با توجه به رشد روزافزون طول زمینه در \lr{LLM}ها (تا ۲M توکن در مدل‌هایی مانند \lr{GPT-5})، بهینه‌سازی حافظه نهان \lr{KV} یکی از حیاتی‌ترین حوزه‌های تحقیقاتی در مهندسی استنتاج خواهد بود.